% =====================================================================
\section{运动学与 Born 截面}\label{sec:dy-born}
% =====================================================================

Drell 与 Yan 在部分子模型中提出了这一过程的基本图像\cite{DrellYan1970}：强子 $A$ 中的一个夸克与强子 $B$ 中的一个反夸克（或反过来）湮灭成虚光子，其余部分子作为旁观者飞走而不参与硬过程。在强子质心系中，两个强子都被 Lorentz 收缩并且时间膨胀，硬湮灭只持续 $\sim1/Q$ 的时间，在这段时间里部分子之间的相互作用来不及改变它们的分布；于是截面应当是对部分子对的非相干求和，每个部分子以它在单个强子中的分布出现。本章要回答的问题是：这一图像在 QCD 的相互作用之下是否仍然成立，在什么意义下成立。最敏感的一点是两个强子的旁观者：它们在硬碰撞之前就已经彼此接近，部分子模型假定它们之间的相互作用不影响湮灭过程。本节先固定运动学与归一化。

\begin{definition}[DY 运动学]\label{def:dy-kin}
考虑 $A(P_A)+B(P_B)\to\gamma^*(q)+X$，$\gamma^*\to\ell^+\ell^-$（只计虚光子交换，一种无质量轻子）。在强子质心系中强子 $A$ 沿 $+z$、$B$ 沿 $-z$ 运动，
$P_A=(P_A^+,\frac{M_A^2}{2P_A^+},\mathbf 0_T)$，$P_B=(\frac{M_B^2}{2P_B^-},P_B^-,\mathbf 0_T)$（光锥分量按\cref{def:prelim-lightfront}）。
\begin{enumerate}[label=(\roman*),leftmargin=2.2em,itemsep=1pt,topsep=2pt]
\item $s=(P_A+P_B)^2$，$Q^2=q^2$，$\tau=Q^2/s$；轻子对的快度 $y=\frac12\ln(q^+/q^-)$。
\item $x_A=\sqrt\tau\,\ee^{y}$，$x_B=\sqrt\tau\,\ee^{-y}$，$0\le\tau\le1$，$|y|\le\frac12\ln\frac1\tau$。
\item 部分子动量 $p_1=\xi_AP_A$、$p_2=\xi_BP_B$（无质量部分子，忽略强子质量），$\hat s=(p_1+p_2)^2=\xi_A\xi_Bs$，$z=Q^2/\hat s$。
\item 部分子截面的归一化：记 $H^\mu$ 为部分子过程 $p_1+p_2\to\gamma^*(\mu)+X$ 的振幅，其中去掉光子传播子与轻子电荷，保留夸克电荷 $e_q$ 与全部 QCD 耦合；$\langle\cdot\rangle$ 表示对初态的自旋与颜色求平均（夸克自旋 $\frac12$、颜色 $\frac1{N_c}$；胶子极化 $\frac{1}{D-2}$、颜色 $\frac1{N_c^2-1}$）并对末态 $X$ 求和；$\dd\Phi_{\gamma^*X}$ 是含 $(2\pi)^D\delta^D(p_1+p_2-q-p_X)$、并固定 $q^2=Q^2$ 的 $D$ 维相空间。令 $\sigma_0\equiv\frac{4\pi\alpha^2}{3N_cQ^2}$，部分子道 $ij$ 的裸系数定义为
\begin{equation}\label{eq:dy-bare}
\Delta^{\rm bare}_{ij}(z,\varepsilon)\equiv\frac{N_c}{2\pi e_q^2(1-\varepsilon)}\int\dd\Phi^{(D)}_{\gamma^*X}\,\big\langle-g_{\mu\nu}H^\mu H^{\nu*}\big\rangle .
\end{equation}
% VF-DY-2
\end{enumerate}
\uses{def:prelim-lightfront}
\end{definition}

定义中的几点说明。由 $q^+=x_1P_A^+$、$q^-=x_2P_B^-$（$x_{1,2}$ 为两个湮灭部分子的动量分数，最低阶时 $q=p_1+p_2$）得 $Q^2=2q^+q^-=x_1x_2\,(2P_A^+P_B^-)$，而 $2P_A^+P_B^-=s\,[1+\order{M^2/s}]$；在 $P_A^+=P_B^-$ 的系中又有 $y=\frac12\ln(x_1/x_2)$。所以最低阶时 $x_1=x_A$、$x_2=x_B$ 正是 (ii) 中由测量量 $(Q^2,y)$ 确定的值，强子质量只以 $x_Ax_B=\tau[1+\order{M^2/s}]$ 的形式出现。\cref{eq:dy-bare}中除以的因子 $(1-\varepsilon)$ 是 $D$ 维 Born 项的自旋迹因子，它使 $\Delta^{(0)}_{q\bar q}=\delta(1-z)$（\cref{prop:dy-born}）；质量因子化的抵消项乘的正是 $D$ 维的 Born 截面，因此只有这样归一，$\MSbar$ 系数函数才不含人为的 $\pm2P_{qq}$（\cref{sec:dy-nlo}）。

\paragraph{主公式。}轻子张量把强子部分与轻子部分分开。对轻子自旋求和，$|M|^2=\frac{e^4}{Q^4}L_{\mu\nu}H^\mu H^{\nu*}$，$L_{\mu\nu}=\Tr[\sla\ell_1\gamma_\mu\sla\ell_2\gamma_\nu]$。在四维中对轻子的两体相空间积分时，Lorentz 协变与 $q^\mu L_{\mu\nu}=0$ 只允许 $\int\dd\Phi_2L_{\mu\nu}=A\,(q_\mu q_\nu-q^2g_{\mu\nu})$；用 $g^{\mu\nu}$ 缩并，左边为 $\int\dd\Phi_2(-8\ell_1\cdot\ell_2)=-\frac{4Q^2}{8\pi}$（无质量两体相空间的体积为 $\frac1{8\pi}$），右边为 $-3Q^2A$，故
\begin{equation}\label{eq:dy-leptontensor}
\int\dd\Phi_2\,L_{\mu\nu}=\frac{1}{6\pi}\big(q_\mu q_\nu-q^2g_{\mu\nu}\big).
\end{equation}
% VF-DY-2
轻子张量对所有部分子道相同，并且在系数函数中消去，所以始终取在四维中。再用 $\frac{\dd^4q}{(2\pi)^4}=\frac{\dd Q^2}{2\pi}\,\frac{\dd^3q}{(2\pi)^32q^0}$、通量 $1/(2\hat s)$ 与流守恒 $q_\mu H^\mu=0$（因而只剩 $-Q^2g_{\mu\nu}$ 一项），得
\[
\frac{\dd\hat\sigma}{\dd Q^2}=\frac{1}{2\hat s}\,\frac{e^4}{Q^4}\,\frac{1}{2\pi}\,\frac{Q^2}{6\pi}\int\dd\Phi_{\gamma^*X}\big\langle-g_{\mu\nu}H^\mu H^{\nu*}\big\rangle ,
\]
以 $e^2=4\pi\alpha$ 代入即
\begin{equation}\label{eq:dy-master}
\hat s\,\frac{\dd\hat\sigma}{\dd Q^2}=\frac{2\alpha^2}{3Q^2}\int\dd\Phi_{\gamma^*X}\,\big\langle-g_{\mu\nu}H^\mu H^{\nu*}\big\rangle ,\qquad
\Delta^{\rm bare}_{ij}=\frac{\hat s\,\dd\hat\sigma_{ij}/\dd Q^2}{\sigma_0\,e_q^2\,(1-\varepsilon)} .
\end{equation}
% VF-DY-2

\begin{proposition}[Born 截面 \tagR]\label{prop:dy-born}
（$R_{\mathrm{PT}}$：在重整化微扰论中对红外正规化的部分子态严格。）本命题是最低阶（$\as^0$）的陈述。
\begin{enumerate}[label=(\alph*),leftmargin=2.2em,itemsep=1pt,topsep=2pt]
\item 对 $q(p_1)\bar q(p_2)\to\gamma^*$，$\hat s\,\dd\hat\sigma^{(0)}/\dd Q^2=\sigma_0\,e_q^2\,(1-\varepsilon)\,\delta(1-z)$，即 $\Delta^{(0)}_{q\bar q}=\delta(1-z)$；在 $D=4$ 中对 $Q^2$ 积分给出 $\hat\sigma(q\bar q\to\ell^+\ell^-)=\dfrac{4\pi\alpha^2e_q^2}{3N_c\hat s}$。
\item 若强子截面具有\cref{def:intro-factorization}的因子化形式，则其最低阶为
\begin{align}
\frac{\dd\sigma}{\dd Q^2}&=\frac{4\pi\alpha^2}{3N_cQ^2s}\sum_qe_q^2\int_0^1\!\dd x_1\,\dd x_2\,\delta(x_1x_2-\tau)\,\big[q(x_1)\bar q(x_2)+\bar q(x_1)q(x_2)\big],\label{eq:dy-LO}\\
\frac{\dd^2\sigma}{\dd Q^2\,\dd y}&=\frac{4\pi\alpha^2}{3N_cQ^2s}\sum_qe_q^2\,\big[q(x_A)\bar q(x_B)+\bar q(x_A)q(x_B)\big],\label{eq:dy-LOy}
\end{align}
% VF-DY-1
其中 $q(x)=f_{q/A}(x)$（\cref{def:pdf-quark}），$\bar q(x)=-f_{q/A}(-x)$（\cref{def:pdf-antiquark}），第一个因子属于强子 $A$、第二个属于 $B$；从 $(x_1,x_2)$ 到 $(x_1x_2,y)$ 的变量变换的 Jacobi 行列式等于 $1$。
\item （Drell--Yan 标度律）$Q^4\,\dd\sigma/\dd Q^2=\frac{4\pi\alpha^2}{3N_c}\,\tau\sum_qe_q^2\int_\tau^1\frac{\dd x}{x}\big[q(x)\,\bar q(\tfrac\tau x)+\bar q(x)\,q(\tfrac\tau x)\big]$，在不计部分子分布的标度依赖时只是 $\tau$ 的函数。
\end{enumerate}
（(a) 只涉及部分子态；(b)(c) 是在因子化形式之下的代数推论，不涉及强子态的非微扰性质。）
\uses{def:dy-kin,def:pdf-quark,def:pdf-antiquark,def:intro-factorization}
\end{proposition}

\begin{proof}
(a) Born 振幅为 $H^\mu=e_q\,\bar v(p_2)\gamma^\mu u(p_1)\,\delta_{ij}$（$i,j$ 为色指标）。在 $D$ 维中 $\gamma^\mu\sla a\gamma_\mu=-(D-2)\sla a$，故对自旋求和
$-g_{\mu\nu}\sum H^\mu H^{\nu*}=-e_q^2\TrD[\sla p_2\gamma^\mu\sla p_1\gamma_\mu]=(D-2)\,e_q^2\,\TrD[\sla p_2\sla p_1]=(D-2)\,e_q^2\,2\hat s$。
颜色求和 $\sum_{ij}\delta_{ij}\delta_{ij}=N_c$，初态平均 $\frac14\cdot\frac1{N_c^2}$，于是 $\langle-gHH^*\rangle=e_q^2(1-\varepsilon)\hat s/N_c$（$D-2=2(1-\varepsilon)$）。$2\to1$ 的相空间为 $\int\dd\Phi_{\gamma^*}=2\pi\,\delta(\hat s-Q^2)$（在任意 $D$ 中成立）。代入\cref{eq:dy-master}，并用 $\delta(\hat s-Q^2)=\delta(1-z)/\hat s$，
$\hat s\,\dd\hat\sigma^{(0)}/\dd Q^2=\frac{2\alpha^2}{3Q^2}\cdot2\pi\cdot\frac{e_q^2(1-\varepsilon)}{N_c}\,\delta(1-z)=\sigma_0e_q^2(1-\varepsilon)\delta(1-z)$，再由\cref{eq:dy-bare}得 $\Delta^{(0)}_{q\bar q}=\delta(1-z)$。$D=4$ 时 $\dd\hat\sigma^{(0)}/\dd Q^2=\sigma_0(Q^2)\,e_q^2\,\delta(\hat s-Q^2)$，对 $Q^2$ 积分即得 $\sigma_0(\hat s)\,e_q^2=4\pi\alpha^2e_q^2/(3N_c\hat s)$。

(b) 因子化形式的最低阶就是部分子截面对两个强子的部分子分布的卷积：$q$ 来自 $A$、$\bar q$ 来自 $B$ 与反过来两种排列，
$\frac{\dd\sigma}{\dd Q^2}=\sum_qe_q^2\int\dd x_1\dd x_2\,[q(x_1)\bar q(x_2)+\bar q(x_1)q(x_2)]\,\sigma_0\,\delta(x_1x_2s-Q^2)$，而 $\delta(x_1x_2s-Q^2)=\delta(x_1x_2-\tau)/s$，即\cref{eq:dy-LO}。对\cref{eq:dy-LOy}，把 $(x_1,x_2)$ 换成 $(u,y)=(x_1x_2,\frac12\ln\frac{x_1}{x_2})$，其 Jacobi 行列式
$\big|\partial(u,y)/\partial(x_1,x_2)\big|=\big|x_2\cdot(-\tfrac{1}{2x_2})-x_1\cdot\tfrac{1}{2x_1}\big|=1$，所以 $\int\dd x_1\dd x_2\,\delta(x_1x_2-\tau)\,\delta(y-\frac12\ln\frac{x_1}{x_2})F(x_1,x_2)=F(\sqrt\tau\ee^{y},\sqrt\tau\ee^{-y})$。

(c) 在\cref{eq:dy-LO}中先对 $x_2$ 积分，$\delta(x_1x_2-\tau)=\delta(x_2-\tau/x_1)/x_1$，乘以 $Q^4=\tau sQ^2$ 即得所述形式；右边除了部分子分布之外只含 $\tau$。
\end{proof}

\cref{prop:dy-born}(b)(c) 只断言从部分子截面到强子截面的代数步骤；“强子截面确实具有因子化形式”本身就是因子化定理在最低阶的内容，其地位是\cref{thm:dy-factorization}的地位（全阶论证，\tagA）。在 $D=4$ 中，用显式旋量计算还可以得到轻子的角分布 $\propto1+\cos^2\theta$（$\theta$ 为轻子对静止系中轻子相对于碰撞轴的极角），这是自旋 $\frac12$ 的部分子湮灭的特征。标度律 (c) 在 $\order{\as}$ 被对数地破坏：系数函数含 $\ln(Q^2/\mu_F^2)$，与部分子分布的演化合在一起才使截面与 $\mu_F$ 无关（\cref{prop:dy-scale}）。
